\section{Making Latin Scripture}\label{sec:jerome-scripture}

\subsection{A Bible before ``the Vulgate''}

Latin Christians already possessed Scripture. Their biblical books circulated
in forms now gathered under the convenient label Old Latin, with revision,
mixture, and regional variation across manuscripts. Jerome did not confront an
empty page or receive a commission to create a complete Bible in one operation.
He entered a living textual tradition whose familiar wording had liturgical,
doctrinal, and emotional weight.

The later word \emph{Vulgate} can therefore mislead in two directions. It may
suggest that Jerome produced every included book, and that one authorized codex
left his desk ready for universal adoption. Neither is true. The collection
that eventually predominated combined Jerome's work with Latin texts made or
revised by others, and it reached relative stability through centuries of
copying, correction, mixture, and use. In Jerome's own day \emph{editio
vulgata}, the ``common edition,'' could refer to the Septuagint-based Old
Testament rather than to a bound Jerome Bible.

\subsection{The Roman Gospel revision}

The first securely bounded project was a revision, not a fresh New Testament
translation. In the preface addressed to Damasus in 383, Jerome describes the
pope's request that divergent Latin copies be compared with Greek witnesses.
He promises the four Gospels, arranged Matthew, Mark, Luke, and John, and says
he changed passages that appeared to alter the sense while leaving much of the
customary Latin intact.

H. A. G. Houghton's manuscript-based analysis confirms the limited scale and
scope: Jerome revised the Gospels in 382--384, intervened unevenly, and appears
to have altered Matthew more heavily than John. His later citations do not
support the old attribution of the Vulgate Acts, Catholic Epistles, Pauline
letters, and Apocalypse to him. At places in his Pauline commentaries he even
criticizes a Latin reading that remains in the later Vulgate form---unlikely
behavior if that form were his own revision. The responsible summary is thus
``Jerome revised the Latin Gospels,'' not ``Damasus commissioned Jerome to
translate the entire New Testament.''

Jerome also worked on Psalters, but their history is layered. A Roman revision,
the widely received Gallican Psalter derived from comparison with Origen's
Hexaplar Greek, and a translation beside the Hebrew are associated with him in
different ways. Attribution and textual descent require manuscript-level
qualification. No single phrase such as ``Jerome's Psalms'' identifies every
Latin psalter later copied under his name.

\subsection{From Greek comparison to Hebrew sources}

At Bethlehem Jerome increasingly translated Old Testament books from Hebrew.
He invoked \emph{Hebraica veritas}: when Greek and Latin witnesses diverged,
the Hebrew source could adjudicate and disclose a wording obscured downstream.
This turn was neither an abandonment of Greek nor a discovery made without
predecessors. Jerome used the Septuagint, other Greek versions, Origen's
Hexapla, existing Latin texts, place-name and name lists, and the exegetical
traditions available around Palestine.

He also learned from Jewish teachers. Letter 125.12 remembers the difficulty of
acquiring Hebrew sounds from a converted Jewish instructor; other prefaces and
commentaries acknowledge consultation with Hebrews. Modern assessment ranges
from admiration of his unusual practical competence to debate about fluency,
dependence on Greek tools, and rhetorical exaggeration. The surviving work
demonstrates serious and sustained Hebrew engagement. It does not demonstrate
that he worked without oral teachers, lexical aids, or uncertainty.

\subsection{What Jerome did and did not translate}

Across roughly 390--405 Jerome produced Hebrew-based Latin versions of most of
the books that form the Catholic Old Testament, along with prefaces explaining
choices and defending the enterprise. The exact project chronology varies by
book. Tobit and Judith were rendered through what he calls ``Chaldean''
(Aramaic), with an intermediary and under haste, rather than by the same method
as a sustained Hebrew-based prophetic translation. Wisdom, Sirach, Baruch, and
the books of Maccabees in the later Vulgate do not become Jerome translations
merely because a medieval Bible placed them beside his work.

The biblical prefaces are therefore part of the biography. They disclose
patrons, assistants, complaints, source choices, speed, uncertainty, and
Jerome's campaign for authority. They also document a canon problem. In the
``helmeted prologue'' to Kings he groups the Hebrew books and places writings
outside that list among apocrypha for purposes of doctrinal confirmation. Yet
Christian use of Wisdom, Sirach, Tobit, Judith, and Maccabees did not disappear,
and Jerome himself translated Tobit and Judith. The Catholic canon later
defined at Trent cannot simply be replaced by Jerome's personal classification.

\subsection{A translator's stated method}

Letter 57.5, written in defense of a disputed translation of Epiphanius,
famously says that Jerome ordinarily renders sense for sense rather than word
for word, while treating Scripture differently because even word order may
matter. The statement is contextual, not a timeless algorithm. His Gospel
preface emphasizes restraint; Hebrew-based prefaces defend return to source;
individual translations range from close formal correspondence to explanatory
rendering. Translation for Jerome was judgment exercised under theological,
textual, and rhetorical pressure.

His commentaries display another dimension. He often offers several inherited
interpretations so that a trained reader can weigh them. This openness did not
make him neutral: he selected, abbreviated, attacked, translated predecessors,
and adapted inherited interpretations. Attribution must distinguish Jerome's
philological decision, his translation of a predecessor, and his adaptation of
that predecessor. No claim about reconstructing a lost Greek text is made
without a directly inspected specialist locus.

\subsection{From contested versions to ecclesial monument}

Readers did not immediately accept Jerome's changes. Augustine worried that a
Hebrew-based wording read publicly would unsettle congregations formed by the
Septuagint; their exchange over Jonah's plant makes the pastoral problem
concrete (Augustine Letter 71; Jerome Letter 75, in the dual-numbered
correspondence). Manuscript mixture continued long after both men died.

Over centuries the collection came to shape western theology, prayer,
education, literature, and law. The Council of Trent received the old and
widely used Latin Vulgate as authentic for public ecclesial use amid sixteenth-
century controversy. As Pope Francis's \emph{Scripturae Sacrae affectus}
explains, that act did not forbid original-language study. The \emph{Nova
Vulgata}, promulgated in 1979 after modern revision, is the Catholic Church's
current typical Latin edition; it is not a claim that every one of its readings
was penned by Jerome.

\evidenceaxis{Jerome was the decisive individual contributor to the later
Vulgate, but not the translator of its whole Bible.}
{Gospel preface to Damasus; Old Testament prefaces; \emph{De viris
illustribus} 135; Houghton, \emph{The Latin New Testament}, pp.~31--35;
Francis, \emph{Scripturae Sacrae affectus}.}
{Jerome's broad self-summary must be tested book by book. The later collection
contains other translators' work and a transmission history no author
controlled.}
